%%%%%%%%%%%%%%%%%%%%%%%%%%%%%%%%%%%%%%%%%%%%%%%%%%%%%%%%%%%%%%%%%%%%%%%
%
% Modelo de preenchimento de informações do cabeçalho
%
% Autor: Vitor G. P. Curtarelli
%
% Data: 10/03/2026
%
%%%%%%%%%%%%%%%%%%%%%%%%%%%%%%%%%%%%%%%%%%%%%%%%%%%%%%%%%%%%%%%%%%%%%%%

\titulo[Developments with TTT and SSF estimation]{Developments with Travel-Time Tomography and Sound Speed Field estimation in water}  % Título do relatório % Argumento em [] (opcional) é o título curto, a ser usado hos headers (vide começo da pág. 2)
\subtitulo{With SVD-based inversion, and context-aware modeling}  % Sub-título do relatório, se houver
\autor{Vitor Curtarelli}  % Autor(es) do relatório
% Vários autores podem ser adicionados em um único comando \autor, ou o comando pode ser chamado várias vezes
%O número de autores no cabeçalho (Autor/Autores) detecta e atualiza automaticamente quando houver 2+ autores
\revisor{Stephan Paul, Lucas C. Lobato} % Revisor(es) do relatório
% Uso idêntico ao comando \autor
\entregavel{3.1.1}{Revisar e organizar SoA de técnicas e algoritmos para inversão do SSF na coluna d'água}  % Número da atividade, como consta na Wiki do projeto
\resumo{This text aims to be an investigative research on SVD-based inversion for Sound Speed field estimation in water through the Travel-Time Tomography technique. It also contains other advancements and applications for the standard model (previously described), along with results and comparisons. It is a continuation of the aforementioned ``Literature Review and standard model description'' document, and assumes its reading.}  % Um breve resumo do que o relatório trata